\par\Needspace{6\baselineskip}\medskip
\begin{center}\hypertarget{native-174-article128}{}\pdfbookmark[2]{128.}{article-128}128.\end{center}
Let us denote the three radius vectors, found as in Articles 125 and 126, by $r,r',r''$. Let the
heliocentric angular motion in the orbit from the second position to the \SourceLeaf{175}third be $2f$, from the
first to the third $2f'$, and from the first to the second $2f''$, so that $f'=f+f''$. Further, let
\[
r'r''\sin2f=n,\qquad rr''\sin2f'=n',\qquad rr'\sin2f''=n'';
\]
and denote the products obtained by multiplying the constant $k$ (Article 2) by the time intervals
from the second
observation to the third, from the first to the third, and from the first to the second respectively by $\theta,\theta',\theta''$. Begin the two calculations of the elements, as in the preceding article, first
from $r,r',f''$ and $\theta''$, and then from $r',r'',f,\theta$. In both calculations, stop short of the elements themselves
as soon as you have found the quantity that expresses the ratio of the elliptic sector to the
triangle and was denoted above (Article 91) by $y$ or $-Y$. Let the value of this quantity be $\eta''$
in the first calculation and $\eta$ in the second. Formula 18 in Article 95 thus gives two values for
the semiparameter $p$
\[
\BareRoot{p}=\frac{\eta''n''}{\theta''}\qquad\text{and}\qquad
\BareRoot{p}=\frac{\eta n}{\theta}\mathpunct{\cdot}
\]
Article 82 further gives a third value
\[
p=\frac{4rr'r''\sin f\sin f'\sin f''}{n-n'+n''}\mathpunct{\raisebox{.45ex}{,}}
\]
and all three values should clearly agree if the true values of $x,y$ were taken from the outset. We
should therefore have
\[
\frac{\theta''}{\theta}=\frac{\eta''n''}{\eta n}
\]
\[
\begin{aligned}
n-n'+n''&=\frac{4\theta\theta''rr'r''\sin f\sin f'\sin f''}{\eta\eta''nn''}\\[5pt]
&=\frac{n'\theta\theta''}{2\eta\eta''rr'r''\cos f\cos f'\cos f''}\mathpunct{\cdot}
\end{aligned}
\]
Unless these equations already happen to be satisfied in the first calculation, we may therefore set
\[
\begin{aligned}
X&=\log\frac{\eta n\theta''}{\eta''n''\theta}\\[5pt]
Y&=n-n'+n''-\frac{n'\theta\theta''}{2\eta\eta''rr'r''\cos f\cos f'\cos f''}\mathpunct{\cdot}
\end{aligned}
\]

This method is as generally applicable as the second method, explained in Article 125. The great
advantage of this fifth method is that the initial hypotheses require the calculation to proceed
only about halfway toward determining the elements themselves. Moreover, as soon as the calculation
has reached the point where we can foresee that the new hypothesis will differ from the truth by no
appreciable amount, it will suffice to determine the elements themselves in that hypothesis from
just $r,r',f'',\theta''$, from $r',r'',f,\theta$, or, preferably, from $r,r'',f',\theta'$.
